\subsection{Chronological orientation}\label{sec:peter-orientation}

No source gives Peter's birth year, age, or a continuous itinerary. The
biographical backbone is instead a sequence of overlapping windows. A
Galilean fisherman enters the public story during Jesus' ministry. He becomes
a leading member of the Twelve, fails during the Passion, and is remembered
as a primary witness of the risen Christ. He then acts publicly in Jerusalem,
participates in the mission beyond it, receives Gentiles, and negotiates the
new movement's relation to Israel and the nations. Paul's letters locate him
in Jerusalem and Antioch. Early Roman Christian memory places his final
witness at Rome. The exact routes between those points are largely unknown.

Jesus' public ministry and crucifixion are commonly placed around
\textsc{ad} 27--30 or 30--33. On either scheme, Peter's call belongs in the
late 20s or early 30s. Paul's first post-call visit to Cephas came, according
to Galatians 1:18, after ``three years''; the meaning and inclusive counting
of that interval are debated. A further Jerusalem meeting ``after fourteen
years'' (Gal 2:1) is usually placed in the later 40s. The incident at Antioch
followed some encounter about table fellowship and Gentile inclusion, but its
precise relation to Acts 10--15 is disputed. Peter then vanishes from Acts as
an active character after chapter 12, apart from his speech in chapter 15.

The Neronian persecution after the fire of Rome in July 64 supplies the most
plausible setting for Peter's death, but the sources do not yield a day or
even a universally certain year. A range of 64--68 is more responsible than a
single date. The fourth-century \emph{Depositio martyrum} and current 29 June
solemnity establish ancient and present liturgical commemoration, not a
surviving civil death certificate or proved bodily translation. The familiar years 42
and 67, used to make a twenty-five-year Roman episcopate, belong to a later
chronological construction and fit awkwardly with the mid-century evidence for
Peter in the East.

\subsection{The principal horizons}

\begin{longtable}{>{\bfseries\raggedright\arraybackslash}p{0.18\linewidth}
                  >{\raggedright\arraybackslash}p{0.31\linewidth}
                  >{\raggedright\arraybackslash}p{0.43\linewidth}}
\toprule
\textbf{Horizon} & \textbf{Principal witness} & \textbf{What it can establish}\\
\midrule
\endhead
Late 20s--early 30s & Gospels; 1 Cor 15:3--8 & Peter's place among Jesus' disciples, his failure in the Passion, and his primacy among named Resurrection witnesses.\\
30s--40s & Acts 1--12; Gal 1:18--24 & Public leadership in Jerusalem and Judea, mission beyond Jerusalem, and Paul's fifteen-day visit with Cephas.\\
Later 40s--50s & Gal 2; Acts 15; 1 Cor 1 and 9 & Peter among the recognized pillars, mission associated with the circumcised, the Antioch conflict, a mobile ministry, and an accompanied or married apostle.\\
64--68 & 1 Clem. 5; Tacitus, \emph{Ann.} 15.44; later Roman witnesses & A probable martyrdom in Nero's Rome. The sources converge more strongly on place and martyr status than on date and manner.\\
Late first--third century & 1 Peter; Ignatius; Irenaeus; Gaius; Tertullian; Origen as preserved by Eusebius & Roman association, the consolidation of apostolic memory, a Vatican memorial, and progressively more detailed accounts of death.\\
Fourth century onward & Calendars, basilicas, liturgies, apocryphal reception, medieval localization & Public cult and material commemoration; also the growth of details that require separate historical grading.\\
\bottomrule
\end{longtable}

The horizons overlap. First Peter, for example, is a canonical text that
locates its sender in ``Babylon'' and transmits a Petrine voice, but its date
and direct authorship are disputed. The Vatican memorial is material evidence
for what Roman Christians of the second century believed, not an independent
biometric identification. First Clement is earlier than the surviving
monument but frustratingly terse. Responsible chronology allows those pieces
to reinforce one another without forcing any one piece to say more than it
does.

\subsection{What remains unknown}

Peter's childhood, education, age at call, appearance, exact economic status,
and most of his travel cannot be recovered. We do not know the name of his
wife, whether she accompanied every journey, or when either spouse died. We
cannot construct a diary between the Antioch incident and Rome. We cannot
demonstrate that Peter served as the sole residential administrator of Rome
for twenty-five uninterrupted years. We cannot establish the date on which he
arrived there or the hour of his death.

These are not trivial gaps to be filled with pious guesses. They define the
scale of the evidence and make the surviving convergences more, not less,
significant. A man whose life is known only in fragments nevertheless stands
at the intersection of independent first-century witnesses and at the center
of a remarkably early Roman remembrance.
